\documentclass[10pt,a4paper]{amsart}
\usepackage[margin=19mm]{geometry}
\usepackage{amsmath,amssymb,mathtools}
\usepackage[scr=rsfs,cal=euler]{mathalfa}
\usepackage{xfrac}
\usepackage[unicode,hidelinks,bookmarks=false]{hyperref}
\newcommand{\ie}{i.e.\ }
\newcommand{\cf}{cf.\ }
\newcommand{\resp}{respectively\ }
\newcommand{\Subsection}{\subsection}
\providecommand{\nobreakdash}{\nobreak\hbox{-}}
\setlength{\emergencystretch}{2em}
\multlinegap0pt
\usepackage[shortlabels]{enumitem}
\usepackage{csquotes}
\usepackage{amssymb}
\usepackage{mathtools}
\usepackage{bbm}
\usepackage[all]{xy}
\usepackage{dsfont}
\usepackage{tikz}
\newtheorem{prop}{Proposition}
\newtheorem{lem}{Lemma}
\newtheorem{thm}{Theorem}
\newcommand{\R}{\mathbb{R}}
\newcommand{\cD}{\mathcal{D}}
\newcommand{\cS}{\mathcal{S}}
\newcommand{\dd}{\mathrm{d}}
\newcommand{\e}{\mathrm{e}}
\newcommand*{\transp}[2][-1mu]{\ensuremath{\mskip1mu\prescript{\smash{\mathrm t\mkern#1}}{}{\mathstrut#2}}}
\title{Values of ternary quadratic forms at integers and the Berry-Tabor conjecture for 3-tori}
\author{Wooyeon Kim, Jens Marklof, Matthew Welsh}
\date{Source: arXiv:2601.03209v1 (CC BY 4.0)}
\begin{document}
\maketitle

\noindent\emph{Source and licence.} Values of ternary quadratic forms at integers and the Berry-Tabor conjecture for 3-tori. Version arXiv:2601.03209v1 (CC BY 4.0), distributed by arXiv under the Creative Commons Attribution 4.0 International licence, http://creativecommons.org/licenses/by/4.0/, as recorded in the arXiv licence metadata for this exact version and shown on the abstract page https://arxiv.org/abs/2601.03209v1. Full licence notice: \texttt{licenses/arxiv-2601.03209-CC-BY-4.0.txt}.\par\smallskip

\noindent\textit{Editorial scope.} Counted unit: the article's prop lemX1 (source0.tex lines 1139--1146). The article's complete original proof of it is delivered with it (source0.tex lines 1148--1235, one \texttt{proof} environment of 88 lines). No mathematics is written, repaired or re-derived: the statement and the proof are the article's own lines, in the article's own notation, and no retained line is substituted or re-flowed.  Retained uncounted as the source-local context the proof needs: lines 706--733; lines 871--887; lines 936--963.  Nothing was omitted from any retained span, so the package carries no omission record.  No retained line contains a citation, so the wrapper carries no bibliography and no bibliography entry.  No retained line contains a figure environment, an image-inclusion command or drawing code, so no image is delivered and the package declares no asset dependency.  Editorial formatting only, in addition: an AMS \texttt{amsart} preamble that restates the article's own declarations and macros used by the retained lines, namely the environments \texttt{prop}, \texttt{lem}, \texttt{thm} on separate counters, together with the standard packages those lines need.  The retained environments print in this document as Proposition 1, Lemma 1, Theorem 1 in order of appearance, whereas the article prints the same items with the numbering of its own full document.  The complete native source of the article as distributed in the arXiv e-print for this exact version is bundled unchanged as \texttt{sources/192212/source0.tex}.
\begin{prop}
  \label{prop:oscillator}
  Let $f\in\cS(\R)$.
  For 
  \begin{equation}
    \label{eq:poscillator}
    p =
    \begin{pmatrix}
      a & b \\
      0 & a^{-1}
    \end{pmatrix}
  \end{equation}
  and any $g \in G$, we have
  \begin{equation}
    \label{eq:poscillator1}
    R_1(p)f(x) = |a|^{\frac{1}{2}} \e( \tfrac{1}{2}  ab x^2) f(ax).
  \end{equation}
  Moreover, for $ g = 
  \begin{pmatrix}
    a & b \\
    c & d
  \end{pmatrix}
  \in G$ with $c \neq 0$, we have
  \begin{equation}
    \label{eq:Fouriertransform}
    R_1(g)f(x) = |c |^{-\frac{1}{2}} \e ( \tfrac{a}{2c}x^2) \int_{\R} f(y) \e( \tfrac{d}{2c}y^2 - \tfrac{1}{c} x y) \dd y.
  \end{equation}
\end{prop}

\begin{lem}
  \label{lem:shortestvector}
  Let $g = (g_j)_j$ and $\gamma = (\gamma_j)_j$.
  Then
  \begin{enumerate}[{\rm (i)}]
  \item $\alpha_l(g) = \min_{\substack{J \subset \{1, \dots, d\} \\ |J| = l}} \prod_{j \in J} \rho(\gamma_j g_j)$,
    
  \item $\rho(\gamma_j g_j) \ll 1$, and
    
  \item if $\rho(g_j) \leq 1$, then $\gamma_j$ has the form $
  \pm \begin{pmatrix}
    1 & * \\
    0 & 1
  \end{pmatrix}
  $.
  \end{enumerate}
\end{lem}

\begin{thm}
  \label{thm:thetaasymptotic}
  Let $F\in\cS(\R^{2d})$ and $A > 0$. Then for $h=(0,y,z)\in H_d$ and
  $$g
  = (
  \begin{pmatrix}
    1 & u_j \\
    0 & 1
  \end{pmatrix}
  \begin{pmatrix}
    v_j^{\frac{1}{2}} & 0 \\
    0 & v_j^{-\frac{1}{2}}
  \end{pmatrix}
  \begin{pmatrix}
    \cos \theta_j & -\sin \theta_j \\
    \sin \theta_j & \cos \theta_j
  \end{pmatrix}
  \in\cD,$$ we have
  \begin{equation}
    \label{eq:thetaasymp}
    \widetilde\Theta_F(h, g) = (v_1 \cdots v_d)^{\frac{1}{2}} F_\theta(0) + O_{F,A}\big((\min v_j)^{-A}\big) 
  \end{equation}
  and, for all $h\in H_d$, $g\in G^g$,
 \begin{equation}
    \label{eq:thetaasymp123}
    |\widetilde\Theta_F(h, g)|  \ll_F \alpha_d(g).
  \end{equation}
\end{thm}

\begin{prop}
Let $d\geq 1$. For $\epsilon>0$ and $A > 1$,
\begin{multline}\label{lemX1}
 \tfrac{1}{L} \int_{|\xi|\leq \epsilon \e^{-t}} \hat{\psi} ( \tfrac{1}{L}\xi ) \widetilde\Theta_F(1, g_0 n(\xi) a(t) g_0^{-1}) \dd \xi \\
= \tfrac{1}{L} \e^{(d-2)t} \hat \psi(0) \int_{\mathbb{R}^{2d}} F(v) \e( \tfrac{1}{2} \xi Q(v) ) \dd v
+ O_{F,\psi, A}(\epsilon^{A} + L^{-2} \e^{(d-4)t} + L^{-1} \epsilon^{-d + 1} \e^t) .
\end{multline}
\end{prop}

\begin{proof}
It will be convenient to write the theta function as 
\begin{equation}
\widetilde\Theta_F(1, g_0 n(\xi) a(t) g_0^{-1}) =
\widetilde\Theta_{F^0}(1, g_0 n(\xi) a(t))
\end{equation}
with $F^0 = R_{d,d}(g_0^{-1}) F$.
We set 
\begin{equation}
  \label{eq:JgXM}
  g = \left(  \begin{pmatrix}   0 & -1 \\
    1 & 0 
  \end{pmatrix}
  \right)_{\!\!1\leq j\leq d}
  g_0  n(\xi) a(t)  = \left(  \begin{pmatrix}   0 & - \e^t x_j^{-\frac{1}{2}} \\
    \e^{-t} x_j^{\frac{1}{2}} & \xi \e^t x_j^{\frac{1}{2}}
  \end{pmatrix}
  \right)_{\!\!1\leq j\leq d},
\end{equation}
and we compute that
\begin{equation}
  \label{eq:gIwasawa}
  g = \left(  \begin{pmatrix}   1 & u_j \\
    0 & 1
  \end{pmatrix}
  \begin{pmatrix}
    v_j^{\frac{1}{2}} & 0 \\
    0 & v_j^{-\frac{1}{2}}
  \end{pmatrix}
  \begin{pmatrix}
    \cos \theta_j & -\sin \theta_j \\
    \sin \theta_j  & \cos \theta_j
  \end{pmatrix}
  \right)_{\!\!1\leq j\leq d},
\end{equation}
with
\begin{equation}
  \label{eq:YRS}
  v_j =  (\e^{-2t} + \xi^2\e^{2t})^{-1}  x_j^{-1},\ \cos \theta_j = \xi \e^t x_j^{\frac{1}{2}}v_j^{\frac{1}{2}},\  \sin \theta_j = \e^{-t} x_j^{\frac{1}{2}}v_j^{\frac{1}{2}}. 
\end{equation}
We observe that if $|\xi| \leq \epsilon \e^{-t}$, then
\begin{equation}
  \label{eq:vnGammang}
  \min_{1\leq j\leq d} v_j \gg \epsilon^{-2}.
\end{equation}
By Lemma \ref{lem:shortestvector}, we may choose $\epsilon$ sufficiently small so that $\alpha_d(g) = \rho_d(g)^{-1}$ and $\alpha_{d-1}(g) \ll \epsilon \alpha_d(g)$. 

By Theorem \ref{thm:thetaasymptotic}, we have
\begin{equation}
  \label{eq:Thetaasymptotic}
  \widetilde\Theta_{F^0}( 1, g_0 n(\xi) a(t ))  = ( v_1 \cdots v_d)^{\frac{1}{2}} F^0_\theta(0) + O_{F, A}(\epsilon^A)
\end{equation}
for any $A > 0$ and from Proposition \ref{prop:oscillator} and \eqref{eq:YRS}, we have
\begin{equation}
  \label{eq:fQ0}
  F^0_\theta(0) = \e^{dt} (v_1 \cdots v_d)^{-\frac{1}{2}} \int_{\mathbb{R}^{2d}} F^0(u,u') \e( \tfrac{1}{2} \xi \e^{2t} (u\transp{u}-u'\transp{u}') ) \dd u \dd u'.
\end{equation}
It follows that
\begin{multline}
  \label{eq:smallxi}
  \tfrac{1}{L} \int_{|\xi|\leq \epsilon \e^{-t}} \hat{\psi} ( \tfrac{1}{L}\xi ) \widetilde\Theta_{F^0}(1, g_0 n(\xi) a(t) ) \dd \xi \\
  = \tfrac{1}{L}\e^{dt} \int_{|\xi| \leq \epsilon \e^{-t}} \hat{\psi}( \tfrac{1}{L}\xi)  \bigg( \int_{\mathbb{R}^{2d}} F^0(u,u') \e( \tfrac{1}{2} \xi \e^{2t} (u\transp{u}-u'\transp{u}') ) \dd u \dd u'\bigg)  \dd \xi + O_{f, \psi, A}(\epsilon^A).
\end{multline}
The inner integral over $\dd u \dd u'$ is $O_{F}( (1 + |\xi| \e^{2t})^{-d})$, so replacing $\hat \psi(\tfrac{1}{L} \xi)$ with $\hat \psi(0)$ creates an error bounded by
\begin{equation}
  \label{eq:psi0replace}
  \ll_{\psi, F} L^{-2} \e^{dt} \int_{|\xi| \leq \epsilon \e^{-t}} \frac{ |\xi|}{(1 + |\xi| \e^{2t})^d} \dd \xi \ll L^{-2} \e^{(d-4)t}. 
\end{equation}
After changing variables $\xi \gets \e^{-2t} \xi$, we have that (\ref{eq:smallxi}) is 
\begin{equation}
\begin{split}
  \label{eq:smallxi1}
&\tfrac{1}{L} \e^{(d-2)t} \hat\psi(0)  \int_{|\xi| \leq \epsilon \e^{t}} \bigg( \int_{\mathbb{R}^{2d}} F^0(u,u') \e( \tfrac{1}{2} \xi (u\transp{u}-u'\transp{u}') ) \dd u \dd u'\bigg)  \dd \xi \\
&  \qquad + O_{F,\psi,A}(\epsilon^{A} + L^{-2} \e^{(d-4)t})\\
& = \tfrac{1}{L} \e^{(d-2)t} \hat \psi(0)  \int_\R \bigg( \int_{\mathbb{R}^{2d}} F^0(u,u') \e( \tfrac{1}{2} \xi (u\transp{u}-u'\transp{u}') ) \dd u \dd u'\bigg)  \dd \xi \\
 & \qquad + O_{F,\psi, A}(\epsilon^{A} + L^{-2} \e^{(d-4)t} + L^{-1} \epsilon^{-d + 1} \e^{-t}).
\end{split}
\end{equation}
Finally, we recall that $F_0=R_{d,d}(g_0^{-1}) F$, which yields after the appropriate variable substitution
\begin{equation}
\begin{split}
 & \int_{\mathbb{R}^{2d}} F^0(u,u') \e( \tfrac{1}{2} \xi (u\transp{u}-u'\transp{u}') ) \dd u \dd u' \\
& = (\det X_0)^{-1} \int_{\mathbb{R}^{2d}} F(u X_0^{-1/2},u' X_0^{-1/2}) \e( \tfrac{1}{2} \xi (u\transp{u}-u'\transp{u}') ) \dd u \dd u' \\
& = \int_{\mathbb{R}^{2d}} F(v) \e( \tfrac{1}{2} \xi Q(v) ) \dd v.
\end{split}
\end{equation}
 
\end{proof}

\end{document}
